
\documentclass[11pt]{article}
\usepackage[letterpaper,margin=0.95in]{geometry}
\usepackage{amsmath,amssymb,amsthm}
\usepackage{graphicx}
\usepackage[authoryear]{natbib}
\usepackage{microtype}
\usepackage[hidelinks]{hyperref}
\usepackage{booktabs}
\usepackage{lineno}
\usepackage{tabularx}
\usepackage{array}
\linespread{1.1}
\graphicspath{{SCOPE_Clean_Supp_Figures/}}

\renewcommand{\thesection}{S\arabic{section}}
\renewcommand{\thesubsection}{S\arabic{section}.\arabic{subsection}}
\renewcommand{\thefigure}{S\arabic{figure}}
\renewcommand{\thetable}{S\arabic{table}}
\numberwithin{equation}{section}

\usepackage{float}
\newtheorem{theorem}{Theorem}
\newtheorem{corollary}{Corollary}
\DeclareMathOperator{\median}{median}
\begin{document}
\title{Supplementary Material for\\
SCOPE: A Family of CDF-Based Wavelet Shrinkage Rules}
\author{Dixon Vimalajeewa$^{1}$, Vijini Lakmini$^{2}$, Malith Premarathna$^{1}$,\\
Fabrizio Ruggeri$^{3}$, and Brani Vidakovic$^{2}$\\[6pt]
{\small $^{1}$Department of Statistics, University of Nebraska, Lincoln, NE 68588, USA}\\
{\small $^{2}$Department of Statistics, Texas A\&M University, College Station, TX 77843, USA}\\
{\small $^{3}$IMATI--CNR, 20133 Milano, Italy}}
\date{November 2026}
\maketitle

\noindent
This supplement accompanies the main article, which develops SCOPE and
N-SCOPE, the calibration of fixed rules, and adaptive SURE selection.
Section~\ref{sec:Bayes} gives the penalized-likelihood and generalized
MAP interpretation of coefficientwise SCOPE,
Section~\ref{supp:proof-scope} proves its structural properties, and
Section~\ref{supp:additional-cdfs} records further CDF constructions.
Section~\ref{sec:s4-calibration} states the calibration rule used for
all reported results, and Section~\ref{sec:supp-simul} gives the complete
numerical results of the root-SNR experiments. Throughout,
$F^*(u)=2F(u)-1$ is the centered CDF of a symmetric unimodal law, and
$\delta(x;\lambda,k)=x|F^*(\lambda x)|^k$, for $\lambda>0$, $k\geq0$.


\section{Penalized-likelihood and MAP interpretation of SCOPE}
\label{sec:Bayes}

A complementary way to understand SCOPE shrinkage is through a Bayesian or penalized likelihood lens. This viewpoint is useful, but it must be stated carefully. For a coefficientwise SCOPE rule satisfying the conditions below, one can construct a symmetric penalty for which the SCOPE estimator is the exact penalized Gaussian estimator. Whether that penalty corresponds to a proper prior density is a separate question, and depends on the tail behavior of the rule. This distinction is important for the prototype rules studied in this paper and will also clarify the role of this section relative to the simulation study.

The analysis below concerns coefficientwise SCOPE, which is the
$r=0$ case of N-SCOPE introduced in Section~2.2 of the main article.
For positive neighborhood radii, attenuation depends on several
observed coefficients, so the scalar inverse construction used here
does not directly apply.

Throughout this section, let
\begin{eqnarray}
\delta(x;\lambda,k) = x\,|F^*(\lambda x)|^k,
\qquad \lambda>0,\;\; k\ge 0,
\label{eq:scope_bayes}
\end{eqnarray}
where $F^*(x)=2F(x)-1$ is the centered version of a continuous symmetric unimodal distribution function $F$.

We assume throughout that for the chosen $F$, $\lambda$, and $k$, the map
$x\mapsto \delta(x;\lambda,k)$ is strictly increasing on ${\mathbb R}$.
From Section~2 of the main article, $\delta$ is continuous, odd, sign preserving, and satisfies
\begin{eqnarray}
|\delta(x;\lambda,k)| &\leq& |x|, \qquad x\in{\mathbb R},\\
\frac{\delta(x;\lambda,k)}{x} &\longrightarrow& 1,
\qquad |x|\to\infty .
\end{eqnarray}


Consider the Gaussian location model
\begin{eqnarray}
X \mid \theta \sim N(\theta,\sigma^2),
\qquad \sigma^2>0.
\label{eq:gauss_loc}
\end{eqnarray}
Since $\delta$ is continuous, odd, unbounded in both directions, and strictly increasing, it admits a continuous odd inverse
\begin{eqnarray}
h(\theta) = \delta^{-1}(\theta).
\label{eq:h_inverse}
\end{eqnarray}
Because $|\delta(x)| \le |x|$, we have
\begin{eqnarray}
h(\theta) \ge \theta
\qquad \mbox{for } \theta \ge 0.
\label{eq:h_ge_theta}
\end{eqnarray}

\medskip
\noindent
\begin{theorem}
\label{th:two}
{\bf (Generalized MAP or penalty representation of SCOPE)}
Let $\delta$ be a SCOPE rule of the form
$$
\delta(x)=x |F^*(\lambda x)|^k, \qquad \lambda>0,\quad k\geq 0,
$$
and assume that $\delta$ is continuous, odd, strictly increasing on ${\mathbb R}$, and satisfies
$$
|\delta(x)|\leq |x|,\qquad x\in{\mathbb R},
$$
together with
$$
\lim_{x\to\infty}\delta(x)=\infty,\qquad
\lim_{x\to-\infty}\delta(x)=-\infty.
$$
Then $\delta$ is a one-to-one and onto map from ${\mathbb R}$ to ${\mathbb R}$, and hence it has a continuous, odd, strictly increasing inverse $h=\delta^{-1}$.
Define
\begin{equation}
\label{eq:scope_penalty}
\Psi(\theta)
=
\frac{1}{\sigma^2}
\int_0^\theta \{h(u)-u\}\,du .
\end{equation}
Then the following statements hold.

First, $\Psi$ is even and nondecreasing in $|\theta|$.

Second, for each observed $x\in{\mathbb R}$, the criterion
$$
Q_x(\theta)=\frac{(x-\theta)^2}{2\sigma^2}+\Psi(\theta)
$$
has the unique global minimizer
$$
\widehat{\theta}(x)=\delta(x).
$$
Equivalently, under these assumptions, the SCOPE rule is an exact penalized Gaussian estimator. It may also be interpreted as a generalized MAP estimator under the formal prior
\begin{eqnarray}
\pi(\theta)\propto \exp\{-\Psi(\theta)\},
\label{eq:formal_prior_scope}
\end{eqnarray}
which may be proper or improper.
\end{theorem}

\medskip
\noindent
\emph{Proof.}
Since $\delta$ is odd and strictly increasing, its inverse $h$ is also odd and strictly increasing. Hence $h(u)-u$ is odd, so the integral in \eqref{eq:scope_penalty} defines an even function $\Psi$.

Next, \eqref{eq:h_ge_theta} implies that $h(\theta)-\theta \ge 0$ for $\theta \ge 0$, hence
\begin{eqnarray}
\Psi'(\theta) = \frac{h(\theta)-\theta}{\sigma^2} \ge 0
\qquad \mbox{for } \theta>0.
\end{eqnarray}
By oddness of $h(\theta)-\theta$, it follows that $\Psi'(\theta) \le 0$ for $\theta<0$. Therefore $\Psi$ is nondecreasing in $|\theta|$.

Now fix $x \in \mathbb{R}$. Differentiating $Q_x(\theta)$ gives
\begin{eqnarray}
Q_x'(\theta)
&=&
\frac{\theta-x}{\sigma^2} + \Psi'(\theta)
\nonumber \\
&=&
\frac{\theta-x}{\sigma^2} + \frac{h(\theta)-\theta}{\sigma^2}
\nonumber \\
&=&
\frac{h(\theta)-x}{\sigma^2}.
\label{eq:Qprime_scope}
\end{eqnarray}
Because $h$ is strictly increasing, $Q_x'(\theta)<0$ when $\theta<\delta(x)$, $Q_x'(\theta)=0$ when $\theta=\delta(x)$, and $Q_x'(\theta)>0$ when $\theta>\delta(x)$. Therefore $Q_x(\theta)$ decreases up to $\theta=\delta(x)$ and increases afterward. Hence $\delta(x)$ is the unique global minimizer of $Q_x(\theta)$.
\hfill $\square$

\medskip
\noindent
Theorem~\ref{th:two} establishes that every coefficientwise SCOPE rule satisfying its hypotheses is an exact penalized estimator under an even, nondecreasing in $|\theta|$ penalty. What is not automatic is that the formal prior in \eqref{eq:formal_prior_scope} is integrable.

\medskip
\noindent
\begin{corollary}
Under the assumptions of Theorem~\ref{th:two}, the formal prior \eqref{eq:formal_prior_scope} is symmetric and unimodal. It is a proper prior density if and only if
\begin{eqnarray}
\int_{-\infty}^{\infty} \exp\{-\Psi(\theta)\}\,d\theta < \infty .
\label{eq:properness_int}
\end{eqnarray}
In particular, if there exist constants $\theta_0>0$ and $c>0$ such that
\begin{eqnarray}
h(\theta)-\theta \ge c
\qquad \mbox{for all } \theta \ge \theta_0,
\label{eq:proper_cond}
\end{eqnarray}
then the prior is proper.
\end{corollary}

\medskip
\noindent
\emph{Proof.}
Symmetry follows from evenness of $\Psi$. Since $\Psi$ is nondecreasing in $|\theta|$, the function $\exp(-\Psi(\theta))$ is nonincreasing on $[0,\infty)$ and nondecreasing on $(-\infty,0]$, hence unimodal with mode at zero.

Condition \eqref{eq:properness_int} is exactly the integrability requirement for a proper prior density. If \eqref{eq:proper_cond} holds, then for $\theta \ge \theta_0$,
\begin{eqnarray}
\Psi(\theta)
&=&
\Psi(\theta_0)
+
\frac{1}{\sigma^2}\int_{\theta_0}^{\theta} \bigl(h(u)-u\bigr)\,du
\nonumber \\
&\ge&
\Psi(\theta_0) + \frac{c}{\sigma^2}(\theta-\theta_0),
\end{eqnarray}
so $\Psi(\theta)$ grows at least linearly in the tails and $\exp(-\Psi(\theta))$ is integrable.
\hfill $\square$

\medskip
\noindent
The corollary shows that the corresponding coefficientwise SCOPE rule induces an even penalty that is nondecreasing in $|\theta|$,
and hence a formal prior that is symmetric and unimodal when written in
the form (\ref{eq:formal_prior_scope}). A positive lower bound on the inverse gap in the tails is sufficient for propriety. When the gap tends to zero, propriety depends on its accumulated integral in the penalty.

\medskip
\noindent
\textbf{Example 1 (Uniform SCOPE: generalized MAP but not proper MAP).}
For the uniform prototype,
\begin{eqnarray}
\delta(x;\lambda,k)
=
\left\{
\begin{array}{ll}
x\,|\lambda x|^k, & |\lambda x| \le 1, \\[1ex]
x, & |\lambda x| > 1 .
\end{array}
\right.
\label{eq:uniform_scope_here}
\end{eqnarray}
Hence $\delta(x)=x$ for all sufficiently large $|x|$, so its inverse satisfies
\begin{eqnarray}
h(\theta)=\theta
\qquad \mbox{for all sufficiently large } |\theta|.
\end{eqnarray}
Therefore $\Psi'(\theta)=0$ in the tails, which implies that $\Psi(\theta)$ is eventually constant. The formal prior \eqref{eq:formal_prior_scope} is then necessarily improper. Thus the uniform SCOPE rule has an exact penalty representation and a generalized MAP interpretation, but not a proper-prior MAP interpretation.

\medskip
\noindent\textbf{Example 2 (Cauchy SCOPE: a proper prior case)}
For the Cauchy prototype with $k>0$,
\begin{eqnarray}
\delta(x;\lambda,k)
=
x\left|\frac{2}{\pi}\arctan(\lambda x)\right|^k .
\label{eq:cauchy_scope_here}
\end{eqnarray}
As $x \to \infty$,
\begin{eqnarray}
\arctan(\lambda x)
=
\frac{\pi}{2} - \frac{1}{\lambda x} + O(x^{-3}),
\end{eqnarray}
hence
\begin{eqnarray}
\frac{2}{\pi}\arctan(\lambda x)
=
1 - \frac{2}{\pi \lambda x} + O(x^{-3}).
\end{eqnarray}
Therefore
\begin{eqnarray}
\delta(x;\lambda,k)
=
x\left(1 - \frac{2}{\pi \lambda x} + O(x^{-3})\right)^k
=
x - \frac{2k}{\pi\lambda} + O(x^{-1})
\qquad \mbox{as } x \to \infty .
\label{eq:cauchy_tail_delta}
\end{eqnarray}
Inverting \eqref{eq:cauchy_tail_delta} gives
\begin{eqnarray}
h(\theta)
=
\theta + \frac{2k}{\pi\lambda} + O(\theta^{-1})
\qquad \mbox{as } \theta \to \infty .
\end{eqnarray}
For $k>0$, this gives $h(\theta)-\theta \ge c$ for all sufficiently large $\theta$, for some $c>0$, and Corollary~1 shows that the induced prior is proper. Thus the Cauchy SCOPE rule is an exact MAP estimator under a proper symmetric unimodal prior.

\medskip
\noindent
\textbf{Remark.}
For the logistic and normal generators, $x-\delta(x)$ decays exponentially or faster as $x\to\infty$. Since $h(\theta)\sim\theta$ and $h(\theta)-\theta=h(\theta)-\delta(h(\theta))$, the inverse gap is integrable in the tail. The penalty therefore has a finite limit and the corresponding formal prior is improper.

\medskip
\noindent
\medskip
\noindent
\textbf{Neighborhood SCOPE.}
For $r>0$, N-SCOPE defines a multivariate mapping whose attenuation
factors depend on overlapping neighborhoods. Although this mapping
retains the sign-preservation and magnitude-contraction properties
established in Section~2.2 of the main article, its Jacobian need not be
symmetric. Thus, the scalar penalty construction of
Theorem~\ref{th:two} does not establish a generalized MAP or a
proper-prior MAP representation for N-SCOPE. Such a representation,
if available under additional restrictions, requires a separate
multivariate analysis.

These penalty and generalized MAP interpretations clarify the
coefficientwise SCOPE subclass but do not determine its tuning. In the
numerical study, the generator is prespecified, while the adaptive procedures select the exponent and scale multiplier and adaptive N-SCOPE also selects the radius, with scales calibrated as described
in Section~4 of the main article.




\section{Proof of the structural properties of SCOPE}
\label{supp:proof-scope}

We prove Proposition~1 of the main article, under its standing assumption
that $F^*$ is continuous, odd, nondecreasing, and satisfies
$uF^*(u)>0$ for $u\ne0$ and $F^*(u)\to1$ as $u\to\infty$.
By oddness of $F^*$,
\begin{eqnarray*}
\delta(-x;\lambda,k)
&=&-x|F^*(-\lambda x)|^k
=-x|F^*(\lambda x)|^k=-\delta(x;\lambda,k).
\end{eqnarray*}
The attenuation $|F^*(\lambda x)|^k$ belongs to $[0,1]$ and is strictly
positive for $x\ne0$ when $k>0$; thus the rule preserves sign and satisfies
$|\delta(x;\lambda,k)|\leq|x|$. The convention $|F^*(\lambda x)|^0\equiv1$
covers $k=0$.

For $x\geq0$, set $a(x)=[F^*(\lambda x)]^k$. Both $a(x)$ and $x$ are
nonnegative and nondecreasing. If $0\leq x_1\leq x_2$, then
\begin{eqnarray*}
\delta(x_1;\lambda,k)=x_1a(x_1)
\leq x_1a(x_2)\leq x_2a(x_2)=\delta(x_2;\lambda,k).
\end{eqnarray*}
Oddness extends monotonicity to the negative half-line and across zero.
Continuity follows from that of $F^*$, while
\begin{eqnarray*}
\frac{\delta(x;\lambda,k)}{x}
=|F^*(\lambda x)|^k\longrightarrow1,
\qquad |x|\longrightarrow\infty,
\end{eqnarray*}
proves relative tail preservation. \hfill$\square$

\section{Additional CDF constructions}
\label{supp:additional-cdfs}

The main article develops the logistic, uniform, and Cauchy prototypes
and lists the normal, Laplace, and hyperbolic-secant generators. We
record two further constructions not treated in the main text:
an exponential-power continuum and mixtures of centered CDFs. For any
smooth symmetric generator with $F^*(u)=c_Fu+o(u)$, $c_F>0$,
\begin{eqnarray}
\delta(x;\lambda,k)
= c_F^k\lambda^k x|x|^k\{1+o(1)\},
\qquad x\longrightarrow0,\quad k>0.
\label{eq:supp-local-general}
\end{eqnarray}
Thus the local attenuation has order $k+1$ for every fixed regular
generator, while the tail shape supplies an additional design choice.
The following extensions are analytical examples; neither was separately
benchmarked in Experiments I--III.

\subsection{An exponential-power continuum}
\label{supp:exponential-power}

An additional class, not studied in the simulations, is generated by the
exponential-power densities
\begin{eqnarray}
f_\beta(u)=\frac{\beta}{2\Gamma(1/\beta)}
\exp\{-|u|^\beta\},\qquad \beta>0.
\label{eq:supp-exp-power-density}
\end{eqnarray}
Their centered CDFs have the explicit representation
\begin{eqnarray}
F_\beta^*(u)
=\operatorname{sgn}(u)\,
\frac{\gamma(1/\beta,|u|^\beta)}{\Gamma(1/\beta)},
\label{eq:supp-exp-power-cdf}
\end{eqnarray}
where $\gamma(a,z)=\int_0^z t^{a-1}e^{-t}\,dt$ denotes the lower
incomplete gamma function. The corresponding SCOPE estimator is
$\delta_\beta(x;\lambda,k)=x|F_\beta^*(\lambda x)|^k$.
Here $\beta=1$ is Laplace and $\beta=2$ is normal up to scale;
moreover, as $\beta\to\infty$, the generating distribution
approaches the uniform law on $[-1,1]$. Thus a single parameter
interpolates between several qualitatively different attenuation
profiles without changing the SCOPE construction. The local slope is
$c_\beta=\beta/\Gamma(1/\beta)$, so
\eqref{eq:supp-local-general} applies to every fixed $\beta>0$.

\subsection{Finite mixtures and multiscale attenuation}
\label{supp:mixtures}

Further flexibility is obtained by mixing centered CDFs. Let
$F_1,\ldots,F_M$ be symmetric unimodal CDFs with a common center,
let $\pi_m\geq0$ with $\sum_{m=1}^M\pi_m=1$, and put
\begin{eqnarray}
F_{\rm mix}^*(u)=\sum_{m=1}^M\pi_m F_m^*(u),
\qquad
\delta_{\rm mix}(x;\lambda,k)
=x\left|\sum_{m=1}^M\pi_mF_m^*(\lambda x)\right|^k.
\label{eq:supp-mixture-basic}
\end{eqnarray}
Because convex mixtures of centered symmetric unimodal CDFs remain
centered symmetric unimodal CDFs, the resulting rule inherits the
coefficientwise SCOPE structural properties. A more flexible version
permits a distinct scale $\lambda_m>0$ for each component:
\begin{eqnarray}
\delta_{\rm mix}(x;\boldsymbol\pi,\boldsymbol\lambda,k)
=x\left|\sum_{m=1}^M\pi_mF_m^*(\lambda_m x)\right|^k.
\label{eq:supp-mixture-multiscale}
\end{eqnarray}
This is again an ordinary SCOPE rule for the centered CDF
$\sum_m\pi_m F_m^*(\lambda_m\cdot)$, although its shape need not
belong to any one original parametric family. In the representation of
the attenuation profile as a mixture of ramps $\min(|u|/t,1)$
(Theorem~1 of the main article, based on \citealp{Khintchine1938}), mixing generators mixes their mixing
laws: if $G_m$ is the mixing law of $F_m$, the multiscale mixture has
mixing law $\sum_m\pi_mG_m^{(\lambda_m)}$, where $G_m^{(\lambda_m)}$
is the law of $T/\lambda_m$ for $T\sim G_m$.

For differentiable generators with slopes $c_m=(F_m^*)'(0)$,
the mixture has origin slope $\sum_m\pi_m\lambda_mc_m$;
for $u>0$, its upper-tail deficit is the weighted sum of the
component deficits. These relations show how a mixture can control
both the central contraction and the persistence of attenuation
into the tails. A component with slowly decaying tails may dominate
the resulting tail behavior even when its mixture weight is modest.

Figure~\ref{fig:supp-mixture-profiles} illustrates the idea with
normal--normal and normal--logistic mixtures. Both use $k=3$,
$\pi_1=\pi_2=1/2$, and component scales
$\lambda_1=0.5$, $\lambda_2=1$. The logistic scale convention is
$F^*(u)=\tanh(u)$, as in the main article. The figure is
illustrative rather than a numerical denoising experiment.
\begin{figure}[htbp]
\centering
\includegraphics[width=0.95\linewidth]{Figs/SCOPE_Mixture_Profiles.pdf}
\caption{Two illustrative centered-CDF mixture SCOPE rules compared with the
identity mapping: normal--normal (left) and normal--logistic (right).
Both use $k=3$, equal mixture weights, and scales $(0.5,1)$.
The curves were regenerated directly from the defining formulas.}
\label{fig:supp-mixture-profiles}
\end{figure}




\section{Calibration and risk estimation}\label{sec:s4-calibration}

This section records the derivations behind the calibration of Section~4 of the main article and the rule used for every reported result.

\subsection{Null risk of a central coefficient and its inversion}
\label{sec:s4-null}
In the orthonormal Gaussian sequence model $d_{j\ell}=\theta_{j\ell}+\sigma Z_{j\ell}$, the N-SCOPE estimate of a central coefficient is $\widehat\theta_{j\ell}=d_{j\ell}[F^*(\lambda_{j,r}S_{j\ell}(r))]^k$, with $S_{j\ell}(r)^2=m^{-1}\sum_{|h|\le r}d_{j,\ell+h}^2$ and $m=2r+1$. We assume $m\le2^{j_0}$ and periodic indexing within each level, so every window contains $m$ distinct coefficients. For $r>0$ each coefficient has its own moving window and its own attenuation factor; no common factor is applied to a disjoint block.

If all $m$ true coefficients in the window vanish, write $Z_0=Z_{j\ell}$ and $V=\sum_{h\ne0}Z_{j,\ell+h}^2\sim\chi^2_{2r}$, independent of $Z_0$, so that $S_{j\ell}(r)^2/\sigma^2=(Z_0^2+V)/m$. With $U=Z_0^2+V\sim\chi^2_m$, spherical symmetry gives $E(Z_0^2\mid U)=U/m$. With $a=\lambda_{j,r}\sigma$, the standardized retained energy of the central coefficient is therefore
\begin{equation}
R_{F,r,k}(a)=E\!\left[\frac{U}{m}\left\{F^*\!\left(a\sqrt{\frac{U}{m}}\right)\right\}^{2k}\right]
=\int_0^\infty\frac{u}{m}\left\{F^*\!\left(a\sqrt{\frac{u}{m}}\right)\right\}^{2k}
\frac{u^{m/2-1}e^{-u/2}}{2^{m/2}\Gamma(m/2)}\,du .
\label{eq:s4-null-integral}
\end{equation}
This calibration accounts for the central coefficient's own contribution to the window energy. For the six generators and any $k>0$, $R_{F,r,k}$ is continuous and strictly increasing in $a>0$, with $R_{F,r,k}(0+)=0$ and $R_{F,r,k}(\infty)=1$ by dominated convergence, since $0\le F^*\le1$ on $[0,\infty)$ and $E(U/m)=1$. The inverse $a_{F,r,k}(p)=R_{F,r,k}^{-1}(p)$, $0<p<1$, is computed by numerical quadrature of \eqref{eq:s4-null-integral} and a bracketed monotone root search; no observed data enter. For a rescaled generator $F_s(u)=F(u/s)$, $R_{F_s,r,k}(a)=R_{F,r,k}(a/s)$ and $a_{F_s,r,k}(p)=s\,a_{F,r,k}(p)$, so the calibrated rule does not depend on the scale convention of $F$; writing the logistic centered CDF as $\tanh(u)$ rather than $\tanh(u/2)$ changes $a$ but not the shrinkage. Because $R_{F,r,k}$ depends on $r$ through $m$, each radius needs its own inverse.

\subsection{Level schedule and noise budgets}
\label{sec:s4-energy}
Write $B_N=(N/2)\eta(N)$ for the expected retained noise energy per shrunk level, in units of $\sigma^2$, with $\eta(N)$ defined in Section~4 of the main article. The anchor uses the normal generator with $k=5$, whose half-attenuation point is $q_5/\lambda$ with $q_5=\Phi^{-1}\{(1+2^{-1/5})/2\}=1.51627656\ldots$; matching it to $\sigma\sqrt{2\log N}$ \citep{DonohoJohnstone1994} gives $a_N=q_5/\sqrt{2\log N}$ and $\eta(N)=E[Z^2\{2\Phi(a_N|Z|)-1\}^{10}]$. The universal threshold serves only as this reference; it is not applied in the estimator. The experiments use the unrounded target $\eta_j=B_N/2^j$, which equals the schedule $\eta_j=\eta(N)2^{J-1-j}$ of the main article, so that
\begin{equation}
\lambda_{j,r}(F,k)=\frac{1}{\widehat\sigma}\,a_{F,r,k}\!\left(\frac{B_N}{2^j}\right).
\label{eq:s4-direct}
\end{equation}
For known noise variance this gives the same expected total retained energy at each shrunk level, even though overlapping estimates are dependent: under the all-null model every estimated coefficient at level $j$ has the same second moment $\sigma^2R_{F,r,k}(\lambda_{j,r}\sigma)$, and linearity of expectation does not require independence, so the level total is $2^jR_{F,r,k}(\lambda_{j,r}\sigma)\sigma^2=B_N\sigma^2$. Substituting the finest-level MAD gives a plug-in calibration.

The target must satisfy $\eta_j<1$ at every shrunk level; since $2^j$ is smallest at $j_0$, this holds if and only if $B_N<2^{j_0}$. For a requested starting level, increase $j_0$ until this holds, leaving coarser coefficients unchanged. Table~\ref{tab:s4-budgets} gives the budgets for the three sample sizes of the experiments and two larger ones.

\begin{table}[H]
\centering\small
\caption{Unrounded budgets $B_N=(N/2)\eta(N)$, the coarsest target $B_N/8$ that $j_0=3$ would imply, and the smallest feasible $j_0\ge3$. The last two columns give the expected retained pure-noise energy under known $\sigma$, as a percentage of the total noise energy $N\sigma^2$, for the shrunk details and for all output coefficients including the $2^{j_0}$ unshrunk scaling coefficients.}
\label{tab:s4-budgets}
\setlength{\tabcolsep}{5pt}
\begin{tabular}{rccccrr}
\toprule
$N$ & $\eta(N)$ & $B_N$ & $B_N/8$ & $j_0$ & Details (\%) & Total (\%) \\
\midrule
512  & 0.0136991 & 3.5070  & 0.4384 & 3 & 4.11 & 5.67 \\
1024 & 0.0100267 & 5.1337  & 0.6417 & 3 & 3.51 & 4.29 \\
2048 & 0.0074791 & 7.6586  & 0.9573 & 3 & 2.99 & 3.38 \\
4096 & 0.0056728 & 11.6178 & 1.4522 & 4 & 2.27 & 2.66 \\
8192 & 0.0043672 & 17.8879 & 2.2360 & 5 & 1.75 & 2.14 \\
\bottomrule
\end{tabular}
\end{table}

\noindent
The percentages follow from $\{(J-j_0)B_N+2^{j_0}\}/N$; at $N=1024$, each of the seven shrunk levels retains $5.13$ units of standardized noise energy. They describe pure-noise output, not the fraction of noise transmitted when an unknown signal is present, which for nonlinear shrinkage depends on the signal.

\subsection{One target for all generators, radii, and exponents}
\label{sec:s4-common}
The calibration uses the same target $B_N/2^j$ for every generator $F$, radius $r$, and exponent $k$, but the solution $a_{F,r,k}(B_N/2^j)$ depends on all three. This is intentional: a common value of $\lambda$ would not produce common noise suppression when the tail of $F$, the distribution of the local energy, or the exponent changes. Null-risk matching places the candidates on the same noise scale while leaving their behavior at nonzero coefficients free to differ. As $r$ increases, the local root-mean-square of pure noise becomes less variable and the calibrated scale increases; an isolated large coefficient can nevertheless be attenuated more strongly at positive radius, while neighboring nonzero coefficients reinforce one another. Equal null risk therefore implies neither equal risk at nonzero coefficients nor the superiority of any radius or generator.

The exponent need not be an integer; $k>0$ suffices because $0\le F^*\le1$ on $[0,\infty)$. By \eqref{eq:supp-local-general}, $k$ sets the order $k+1$ at which the estimate leaves zero. Since the scale is recalibrated for each $k$, the null-risk target does not by itself select an exponent for nonzero coefficients; fixed SCOPE uses $k=5$ as a prespecified convention.

\subsection{SURE score and its diagonal derivative}
\label{sec:s4-sure}
For a candidate $(k,r,b)$ the scale is $\lambda_{j,r}(F,k,b)=b\,a_{F,r,k}(B_N/2^j)/\widehat\sigma$, so only the global multiplier $b$ moves; level-to-level differences stay fixed by $B_N/2^j$. With $S=S_{j\ell}(r)$, $\lambda=\lambda_{j,r}(F,k,b)$, $u=\lambda S$, and $w=[F^*(u)]^k$, and with the candidate scale held fixed,
\begin{equation}
\frac{\partial S}{\partial d_{j\ell}}=\frac{d_{j\ell}}{mS},\qquad
\frac{\partial(d_{j\ell}w)}{\partial d_{j\ell}}
=[F^*(u)]^k+\frac{k\lambda d_{j\ell}^2}{mS}[F^*(u)]^{k-1}(F^*)'(u),
\label{eq:s4-sure-diagonal}
\end{equation}
which is the derivative used in Section~4 of the main article. It includes the effect of the central coefficient on its own window energy. Overlap creates off-diagonal partial derivatives, but Stein's divergence is the trace of the Jacobian and needs only the diagonal. At $S=0$ the derivative is defined by its limit; for the uniform generator, the derivative almost everywhere suffices under the continuous Gaussian model. In the implementation the MAD estimate $\widehat\sigma=\median_\ell|d_{J-1,\ell}|/0.67448975$ and the candidate scales are held fixed when evaluating \eqref{eq:s4-sure-diagonal}, and the summed score over all shrunk levels is minimized over $k\in\{3,5,7\}$, $b\in\{0.65,0.80,1.00,1.20,1.50\}$, and, for N-SCOPE SURE, $r\in\{0,1,2,3\}$. The selected triple is common to all detail levels of a realization. The criterion is therefore a plug-in score, not an exactly unbiased risk estimate for the procedure with estimated noise and selected parameters. Performance of the selected rule is evaluated on the simulated truths, using the same noisy observations for all competing procedures.

\section{Additional numerical evidence for the simulation study}
\label{sec:supp-simul}

This section gives the numerical details behind Section~5 of the main article.
All quantities below are taken from the completed $M=1000$ output for
$N=512,1024,2048$ and \emph{root} SNR $3,5,7$.
The MATLAB implementation multiplies the clean signal by its
requested ratio of signal standard deviation to noise standard deviation,
with independent noise $N(0,1)$. It does \emph{not} impose power SNR
$N^{-1}\sum_i f(t_i)^2/\sigma^2$. The tables use the same periodic
orthonormal wavelet transforms and leave eight scaling coefficients
unshrunk. MSE is computed in wavelet coordinates and divided by $N$;
Parseval's identity makes it identical to reconstruction MSE.

\subsection{Calibration targets and finite-size conventions}
\label{sec:supp-calib}
The finest-level Gaussian null-risk targets are computed before seeing
the signal. With $J=\log_2N$, the calibration uses
$\eta_j=\eta(N)2^{J-1-j}$ for $j=3,\ldots,J-1$.
The following table makes explicit how close the coarse target is to
one at $N=2048$. The positive-radius constants are obtained separately
for each generator, radius, exponent, and level by numerical inversion
of the same $\chi^2_{2r+1}$ null-risk equation. The central coefficient
alone is attenuated, even though its calibration uses neighborhood energy.

\begin{table}[H]
\centering\small
\caption{Preassigned Gaussian null-risk targets; $j_0=3$ and $J=\log_2 N$. The coarsest target is $\eta_3(N)=2^{J-4}\eta(N)$, computed from the unrounded $\eta(N)$ of Table~\ref{tab:s4-budgets}. These are per-coefficient normalized null risks, not reconstruction MSEs.}
\label{tab:s5-eta}
\setlength{\tabcolsep}{5pt}
\begin{tabular}{rrrrrr}
\toprule
$N$ & $J$ & Finest $j$ & $\eta(N)$ & Coarsest $\eta_3$ & Shrunk levels \\
\midrule
512 & 9 & 8 & 0.01370 & 0.4384 & 6 \\
1024 & 10 & 9 & 0.01003 & 0.6417 & 7 \\
2048 & 11 & 10 & 0.00748 & 0.9573 & 8 \\
\bottomrule
\end{tabular}
\end{table}

\subsection{Experiment I: all sizes and representative generating families}
\label{sec:supp-e1}
Tables~\ref{tab:s5-e1-normal-512}--\ref{tab:s5-e1-normal-2048}
report \emph{all twelve} signal--root-SNR combinations for the Normal
generator at each sample size. Corresponding logistic tables emphasize
that the neighborhood pattern is not specific to a single centered
CDF. Every row compares $r=0,1,2,3$ after independent null-risk
calibration; no radius minimizes an oracle criterion. In particular,
the Blocks penalty from positive radii should not be hidden by an
aggregate score.

\begin{table}[H]
\centering\small
\caption{Experiment I, Normal generator, $N=512$, $k=5$, $M=1000$: AMSE for fixed radii at all four signals and root SNR $3,5,7$.}
\label{tab:s5-e1-normal-512}
\setlength{\tabcolsep}{5pt}
\begin{tabular}{llcccc}
\toprule
Signal & Root SNR & $r=0$ & $r=1$ & $r=2$ & $r=3$ \\
\midrule
Doppler & 3 & 0.2219 & 0.1820 & 0.1735 & 0.1762 \\
Doppler & 5 & 0.2809 & 0.2141 & 0.1952 & 0.1933 \\
Doppler & 7 & 0.3091 & 0.2285 & 0.2093 & 0.2065 \\
HeaviSine & 3 & 0.0997 & 0.0979 & 0.1008 & 0.1029 \\
HeaviSine & 5 & 0.1093 & 0.1120 & 0.1211 & 0.1283 \\
HeaviSine & 7 & 0.1181 & 0.1241 & 0.1367 & 0.1477 \\
Bumps & 3 & 0.4770 & 0.3837 & 0.4469 & 0.4946 \\
Bumps & 5 & 0.5148 & 0.3709 & 0.4409 & 0.5099 \\
Bumps & 7 & 0.5474 & 0.3617 & 0.4299 & 0.5021 \\
Blocks & 3 & 0.2025 & 0.2786 & 0.3245 & 0.3600 \\
Blocks & 5 & 0.1791 & 0.2758 & 0.3305 & 0.3795 \\
Blocks & 7 & 0.1649 & 0.2655 & 0.3292 & 0.3774 \\
\bottomrule
\end{tabular}
\end{table}

\begin{table}[H]
\centering\small
\caption{Experiment I, Normal generator, $N=1024$, $k=5$, $M=1000$: AMSE for fixed radii at all four signals and root SNR $3,5,7$.}
\label{tab:s5-e1-normal-1024}
\setlength{\tabcolsep}{5pt}
\begin{tabular}{llcccc}
\toprule
Signal & Root SNR & $r=0$ & $r=1$ & $r=2$ & $r=3$ \\
\midrule
Doppler & 3 & 0.1357 & 0.1132 & 0.1093 & 0.1112 \\
Doppler & 5 & 0.1610 & 0.1331 & 0.1278 & 0.1290 \\
Doppler & 7 & 0.1804 & 0.1489 & 0.1404 & 0.1398 \\
HeaviSine & 3 & 0.0634 & 0.0650 & 0.0685 & 0.0713 \\
HeaviSine & 5 & 0.0696 & 0.0737 & 0.0799 & 0.0852 \\
HeaviSine & 7 & 0.0772 & 0.0808 & 0.0889 & 0.0957 \\
Bumps & 3 & 0.2828 & 0.2252 & 0.2623 & 0.2964 \\
Bumps & 5 & 0.3126 & 0.2356 & 0.2797 & 0.3218 \\
Bumps & 7 & 0.3179 & 0.2422 & 0.2868 & 0.3311 \\
Blocks & 3 & 0.1588 & 0.2210 & 0.2605 & 0.2904 \\
Blocks & 5 & 0.1240 & 0.2021 & 0.2585 & 0.3021 \\
Blocks & 7 & 0.1054 & 0.1847 & 0.2470 & 0.2914 \\
\bottomrule
\end{tabular}
\end{table}

\begin{table}[H]
\centering\small
\caption{Experiment I, Normal generator, $N=2048$, $k=5$, $M=1000$: AMSE for fixed radii at all four signals and root SNR $3,5,7$.}
\label{tab:s5-e1-normal-2048}
\setlength{\tabcolsep}{5pt}
\begin{tabular}{llcccc}
\toprule
Signal & Root SNR & $r=0$ & $r=1$ & $r=2$ & $r=3$ \\
\midrule
Doppler & 3 & 0.0813 & 0.0723 & 0.0697 & 0.0702 \\
Doppler & 5 & 0.0925 & 0.0828 & 0.0794 & 0.0789 \\
Doppler & 7 & 0.1061 & 0.0897 & 0.0851 & 0.0843 \\
HeaviSine & 3 & 0.0444 & 0.0460 & 0.0486 & 0.0506 \\
HeaviSine & 5 & 0.0503 & 0.0523 & 0.0559 & 0.0591 \\
HeaviSine & 7 & 0.0562 & 0.0576 & 0.0625 & 0.0668 \\
Bumps & 3 & 0.1689 & 0.1393 & 0.1610 & 0.1811 \\
Bumps & 5 & 0.1854 & 0.1541 & 0.1798 & 0.2046 \\
Bumps & 7 & 0.1913 & 0.1614 & 0.1879 & 0.2157 \\
Blocks & 3 & 0.0932 & 0.1309 & 0.1573 & 0.1765 \\
Blocks & 5 & 0.0802 & 0.1259 & 0.1603 & 0.1880 \\
Blocks & 7 & 0.0697 & 0.1177 & 0.1560 & 0.1873 \\
\bottomrule
\end{tabular}
\end{table}

\begin{table}[H]
\centering\small
\caption{Experiment I, Logistic generator, $N=512$, $k=5$, $M=1000$: AMSE for fixed radii at all four signals and root SNR $3,5,7$.}
\label{tab:s5-e1-logistic-512}
\setlength{\tabcolsep}{5pt}
\begin{tabular}{llcccc}
\toprule
Signal & Root SNR & $r=0$ & $r=1$ & $r=2$ & $r=3$ \\
\midrule
Doppler & 3 & 0.2195 & 0.1812 & 0.1732 & 0.1759 \\
Doppler & 5 & 0.2827 & 0.2159 & 0.1968 & 0.1947 \\
Doppler & 7 & 0.3147 & 0.2327 & 0.2114 & 0.2072 \\
HeaviSine & 3 & 0.0975 & 0.0958 & 0.0983 & 0.1001 \\
HeaviSine & 5 & 0.1085 & 0.1108 & 0.1188 & 0.1253 \\
HeaviSine & 7 & 0.1178 & 0.1233 & 0.1351 & 0.1454 \\
Bumps & 3 & 0.4898 & 0.3958 & 0.4599 & 0.5038 \\
Bumps & 5 & 0.5311 & 0.3824 & 0.4549 & 0.5257 \\
Bumps & 7 & 0.5629 & 0.3704 & 0.4404 & 0.5151 \\
Blocks & 3 & 0.2084 & 0.2829 & 0.3279 & 0.3621 \\
Blocks & 5 & 0.1858 & 0.2826 & 0.3370 & 0.3869 \\
Blocks & 7 & 0.1700 & 0.2716 & 0.3338 & 0.3841 \\
\bottomrule
\end{tabular}
\end{table}

\begin{table}[H]
\centering\small
\caption{Experiment I, Logistic generator, $N=1024$, $k=5$, $M=1000$: AMSE for fixed radii at all four signals and root SNR $3,5,7$.}
\label{tab:s5-e1-logistic-1024}
\setlength{\tabcolsep}{5pt}
\begin{tabular}{llcccc}
\toprule
Signal & Root SNR & $r=0$ & $r=1$ & $r=2$ & $r=3$ \\
\midrule
Doppler & 3 & 0.1357 & 0.1127 & 0.1085 & 0.1101 \\
Doppler & 5 & 0.1620 & 0.1335 & 0.1278 & 0.1287 \\
Doppler & 7 & 0.1824 & 0.1496 & 0.1408 & 0.1401 \\
HeaviSine & 3 & 0.0628 & 0.0641 & 0.0672 & 0.0696 \\
HeaviSine & 5 & 0.0692 & 0.0731 & 0.0790 & 0.0839 \\
HeaviSine & 7 & 0.0767 & 0.0804 & 0.0882 & 0.0948 \\
Bumps & 3 & 0.2872 & 0.2290 & 0.2661 & 0.2998 \\
Bumps & 5 & 0.3199 & 0.2393 & 0.2838 & 0.3268 \\
Bumps & 7 & 0.3276 & 0.2449 & 0.2902 & 0.3358 \\
Blocks & 3 & 0.1633 & 0.2252 & 0.2643 & 0.2935 \\
Blocks & 5 & 0.1306 & 0.2096 & 0.2663 & 0.3115 \\
Blocks & 7 & 0.1099 & 0.1899 & 0.2517 & 0.2972 \\
\bottomrule
\end{tabular}
\end{table}

\begin{table}[H]
\centering\small
\caption{Experiment I, Logistic generator, $N=2048$, $k=5$, $M=1000$: AMSE for fixed radii at all four signals and root SNR $3,5,7$.}
\label{tab:s5-e1-logistic-2048}
\setlength{\tabcolsep}{5pt}
\begin{tabular}{llcccc}
\toprule
Signal & Root SNR & $r=0$ & $r=1$ & $r=2$ & $r=3$ \\
\midrule
Doppler & 3 & 0.0815 & 0.0720 & 0.0692 & 0.0695 \\
Doppler & 5 & 0.0927 & 0.0827 & 0.0791 & 0.0785 \\
Doppler & 7 & 0.1060 & 0.0900 & 0.0853 & 0.0843 \\
HeaviSine & 3 & 0.0441 & 0.0456 & 0.0480 & 0.0498 \\
HeaviSine & 5 & 0.0499 & 0.0518 & 0.0554 & 0.0584 \\
HeaviSine & 7 & 0.0557 & 0.0572 & 0.0620 & 0.0662 \\
Bumps & 3 & 0.1703 & 0.1401 & 0.1616 & 0.1817 \\
Bumps & 5 & 0.1893 & 0.1556 & 0.1815 & 0.2067 \\
Bumps & 7 & 0.1963 & 0.1633 & 0.1903 & 0.2185 \\
Blocks & 3 & 0.0957 & 0.1331 & 0.1593 & 0.1783 \\
Blocks & 5 & 0.0835 & 0.1296 & 0.1642 & 0.1922 \\
Blocks & 7 & 0.0725 & 0.1211 & 0.1593 & 0.1908 \\
\bottomrule
\end{tabular}
\end{table}


To summarize all six generators without privileging a single root-SNR
condition, the following tables give the equal-weighted mean over
root SNR $3,5,7$ of the \emph{percentage} gain in AMSE of fixed
$r=1$ relative to fixed $r=0$, evaluated separately for each signal.
A negative percentage means deterioration. Averaging percentages
rather than raw squared errors avoids allowing the highest-signal
conditions to dominate the summary.

\begin{table}[H]
\centering\small
\caption{Experiment I, $N=512$: average AMSE reduction (\%) from fixed $r=1$ versus $r=0$, across root SNR $3,5,7$, for all centered-CDF generators; positive favors N-SCOPE.}
\label{tab:s5-gain-512}
\setlength{\tabcolsep}{5pt}
\begin{tabular}{lrrrr}
\toprule
Generator & Doppler & HeaviSine & Bumps & Blocks \\
\midrule
Logistic & +22.4 & -1.7 & +27.1 & -49.2 \\
Normal & +22.6 & -1.9 & +27.1 & -50.8 \\
Laplace & +21.7 & -1.5 & +26.3 & -47.7 \\
Sech & +22.2 & -1.6 & +26.9 & -48.4 \\
Uniform & +24.1 & -2.2 & +29.2 & -52.7 \\
Cauchy & +17.9 & -1.2 & +14.3 & -29.5 \\
\bottomrule
\end{tabular}
\end{table}

\begin{table}[H]
\centering\small
\caption{Experiment I, $N=1024$: average AMSE reduction (\%) from fixed $r=1$ versus $r=0$, across root SNR $3,5,7$, for all centered-CDF generators; positive favors N-SCOPE.}
\label{tab:s5-gain-1024}
\setlength{\tabcolsep}{5pt}
\begin{tabular}{lrrrr}
\toprule
Generator & Doppler & HeaviSine & Bumps & Blocks \\
\midrule
Logistic & +17.5 & -4.2 & +23.6 & -57.1 \\
Normal & +17.1 & -4.4 & +22.9 & -59.2 \\
Laplace & +17.8 & -4.2 & +24.1 & -55.6 \\
Sech & +17.7 & -4.2 & +23.9 & -56.2 \\
Uniform & +17.0 & -4.9 & +23.1 & -62.8 \\
Cauchy & +18.8 & -3.3 & +16.9 & -32.9 \\
\bottomrule
\end{tabular}
\end{table}

\begin{table}[H]
\centering\small
\caption{Experiment I, $N=2048$: average AMSE reduction (\%) from fixed $r=1$ versus $r=0$, across root SNR $3,5,7$, for all centered-CDF generators; positive favors N-SCOPE.}
\label{tab:s5-gain-2048}
\setlength{\tabcolsep}{5pt}
\begin{tabular}{lrrrr}
\toprule
Generator & Doppler & HeaviSine & Bumps & Blocks \\
\midrule
Logistic & +12.5 & -3.4 & +17.4 & -53.8 \\
Normal & +12.4 & -3.3 & +16.7 & -55.5 \\
Laplace & +12.7 & -3.5 & +18.3 & -52.4 \\
Sech & +12.6 & -3.4 & +17.9 & -53.0 \\
Uniform & +12.9 & -3.2 & +16.6 & -59.2 \\
Cauchy & +16.3 & -3.5 & +17.1 & -35.0 \\
\bottomrule
\end{tabular}
\end{table}


\subsection{Experiment II: normalized risks and parameter selection}
\label{sec:supp-e2}
For each generator and each $N$, we compute the normalized
$Q_N$, defined as the equal-weighted mean of twelve conditionwise AMSE ratios to universal hard thresholding, over the same twelve
signal--root-SNR conditions. The two fixed procedures use $k=5$,
$b=1$, and $r=0$ or $1$; SURE selects $(k,b)$ for ordinary SCOPE
or $(k,r,b)$ for N-SCOPE. The generator itself is not optimized
within a noisy realization. Each following table contains the complete
six-generator comparison, rather than a selected winner.

\begin{table}[H]
\centering\small
\caption{Experiment II, $N=512$: macro-averaged normalized AMSE $Q_N$ (universal-hard baseline $Q_N=1$) over twelve signal--root-SNR conditions.}
\label{tab:s5-e2-q-512}
\setlength{\tabcolsep}{5pt}
\begin{tabular}{lrrrr}
\toprule
Generator & Fixed SCOPE & Fixed N-SCOPE & SCOPE SURE & N-SCOPE SURE \\
\midrule
Logistic & 0.7210 & 0.7467 & 0.6862 & 0.6181 \\
Normal & 0.7105 & 0.7377 & 0.6816 & 0.6106 \\
Laplace & 0.7486 & 0.7766 & 0.6965 & 0.6326 \\
Sech & 0.7312 & 0.7570 & 0.6910 & 0.6247 \\
Uniform & 0.7239 & 0.7444 & 0.7197 & 0.6333 \\
Cauchy & 1.2747 & 1.3352 & 0.9881 & 0.9285 \\
\bottomrule
\end{tabular}
\end{table}

\begin{table}[H]
\centering\small
\caption{Experiment II, $N=1024$: macro-averaged normalized AMSE $Q_N$ (universal-hard baseline $Q_N=1$) over twelve signal--root-SNR conditions.}
\label{tab:s5-e2-q-1024}
\setlength{\tabcolsep}{5pt}
\begin{tabular}{lrrrr}
\toprule
Generator & Fixed SCOPE & Fixed N-SCOPE & SCOPE SURE & N-SCOPE SURE \\
\midrule
Logistic & 0.7534 & 0.8189 & 0.7090 & 0.6425 \\
Normal & 0.7412 & 0.8089 & 0.7005 & 0.6328 \\
Laplace & 0.7843 & 0.8533 & 0.7255 & 0.6610 \\
Sech & 0.7649 & 0.8306 & 0.7159 & 0.6502 \\
Uniform & 0.7461 & 0.8145 & 0.7179 & 0.6406 \\
Cauchy & 1.3527 & 1.4488 & 1.0781 & 1.0099 \\
\bottomrule
\end{tabular}
\end{table}

\begin{table}[H]
\centering\small
\caption{Experiment II, $N=2048$: macro-averaged normalized AMSE $Q_N$ (universal-hard baseline $Q_N=1$) over twelve signal--root-SNR conditions.}
\label{tab:s5-e2-q-2048}
\setlength{\tabcolsep}{5pt}
\begin{tabular}{lrrrr}
\toprule
Generator & Fixed SCOPE & Fixed N-SCOPE & SCOPE SURE & N-SCOPE SURE \\
\midrule
Logistic & 0.8121 & 0.8974 & 0.7441 & 0.6727 \\
Normal & 0.8025 & 0.8891 & 0.7316 & 0.6614 \\
Laplace & 0.8381 & 0.9270 & 0.7631 & 0.6929 \\
Sech & 0.8215 & 0.9074 & 0.7529 & 0.6815 \\
Uniform & 0.8104 & 0.8958 & 0.7378 & 0.6601 \\
Cauchy & 1.3689 & 1.4847 & 1.1559 & 1.0839 \\
\bottomrule
\end{tabular}
\end{table}


Next we report the probability distribution of the \emph{selected}
N-SCOPE SURE radius at each sample size, pooled equally across the
three root-SNR settings. Since $M=1000$ per setting, these are
frequencies over 3000 independent simulated reconstructions for each
signal and generator. The radius is one global parameter for a
reconstruction; it is not a separate block decision.

\begin{table}[H]
\centering\small
\caption{Experiment II, $N=512$: N-SCOPE SURE selected-radius frequencies (\%) pooled over root SNR $3,5,7$. Each row sums to approximately $100\%$.}
\label{tab:s5-radius-512}
\setlength{\tabcolsep}{5pt}
\begin{tabular}{llrrrr}
\toprule
Generator & Signal & $r=0$ & $r=1$ & $r=2$ & $r=3$ \\
\midrule
Normal & Doppler & 0.1 & 7.0 & 35.8 & 57.1 \\
Normal & HeaviSine & 41.5 & 37.3 & 11.4 & 9.9 \\
Normal & Bumps & 3.8 & 96.2 & 0.0 & 0.0 \\
Normal & Blocks & 99.9 & 0.1 & 0.0 & 0.0 \\
Logistic & Doppler & 0.2 & 7.4 & 36.3 & 56.1 \\
Logistic & HeaviSine & 40.9 & 38.1 & 11.0 & 10.0 \\
Logistic & Bumps & 3.5 & 96.5 & 0.0 & 0.0 \\
Logistic & Blocks & 99.9 & 0.1 & 0.0 & 0.0 \\
\bottomrule
\end{tabular}
\end{table}

\begin{table}[H]
\centering\small
\caption{Experiment II, $N=1024$: N-SCOPE SURE selected-radius frequencies (\%) pooled over root SNR $3,5,7$. Each row sums to approximately $100\%$.}
\label{tab:s5-radius-1024}
\setlength{\tabcolsep}{5pt}
\begin{tabular}{llrrrr}
\toprule
Generator & Signal & $r=0$ & $r=1$ & $r=2$ & $r=3$ \\
\midrule
Normal & Doppler & 0.0 & 5.9 & 42.6 & 51.5 \\
Normal & HeaviSine & 54.9 & 34.4 & 7.2 & 3.5 \\
Normal & Bumps & 0.2 & 99.8 & 0.0 & 0.0 \\
Normal & Blocks & 100.0 & 0.0 & 0.0 & 0.0 \\
Logistic & Doppler & 0.0 & 5.6 & 42.0 & 52.4 \\
Logistic & HeaviSine & 55.2 & 34.0 & 7.3 & 3.5 \\
Logistic & Bumps & 0.1 & 99.9 & 0.0 & 0.0 \\
Logistic & Blocks & 100.0 & 0.0 & 0.0 & 0.0 \\
\bottomrule
\end{tabular}
\end{table}

\begin{table}[H]
\centering\small
\caption{Experiment II, $N=2048$: N-SCOPE SURE selected-radius frequencies (\%) pooled over root SNR $3,5,7$. Each row sums to approximately $100\%$.}
\label{tab:s5-radius-2048}
\setlength{\tabcolsep}{5pt}
\begin{tabular}{llrrrr}
\toprule
Generator & Signal & $r=0$ & $r=1$ & $r=2$ & $r=3$ \\
\midrule
Normal & Doppler & 0.1 & 1.7 & 20.3 & 77.9 \\
Normal & HeaviSine & 60.3 & 34.8 & 3.6 & 1.3 \\
Normal & Bumps & 0.2 & 99.8 & 0.0 & 0.0 \\
Normal & Blocks & 100.0 & 0.0 & 0.0 & 0.0 \\
Logistic & Doppler & 0.1 & 1.6 & 19.7 & 78.6 \\
Logistic & HeaviSine & 61.0 & 33.9 & 3.8 & 1.4 \\
Logistic & Bumps & 0.1 & 99.8 & 0.0 & 0.0 \\
Logistic & Blocks & 100.0 & 0.0 & 0.0 & 0.0 \\
\bottomrule
\end{tabular}
\end{table}


The pooled radius frequencies summarize selection across replications and signal strengths. The SURE-selected exponent $k$ and scale multiplier $b$ are available as replication-specific records in the numerical data accompanying this supplement, together with the joint frequencies of the selected triples; these are tuning diagnostics, not posterior probabilities or hypothesis tests.
% CONFIRM (authors): these records are in the public repository; otherwise delete the previous sentence. Because $\widehat\sigma$ is estimated
from the same sample and its derivative is held fixed, the criterion
is plug-in SURE and is not exactly unbiased after selection.

The selection frequencies agree with the AMSEs. In Experiment~II at $N=1024$ and root SNR~$5$, a fixed radius $r=1$ already gives nearly the full gain on Bumps (AMSE $0.2356$ versus $0.2351$ with joint adaptation), and N-SCOPE SURE selects $r=1$ in nearly every realization. For Blocks, where $r=0$ is selected almost always, allowing radius zero recovers the coefficientwise AMSE of $0.1211$, compared with $0.2021$ for fixed $r=1$. Doppler, which mostly selects $r=2$ or $3$, gains further from joint adaptation, whereas coefficientwise SURE performs slightly better on HeaviSine, where the selected radius varies most. Because each signal is analyzed with its own wavelet, these are findings for the tested signal--wavelet combinations.

\subsection{Experiment III: complete implementation-level rankings}
\label{sec:supp-e3}
Experiment III uses fresh Gaussian-noise streams and four frozen
Normal-generator SCOPE procedures alongside twenty literature-based
codes and universal hard/soft references. The following ranking table
includes \emph{every executed code}, with each $Q_N$ calculated
separately from the twelve signal--root-SNR conditions for that $N$.
The HMT, context and EB-block entries designated as proxies are not
validated implementations of their cited published algorithms;
additional implementation-specific qualifications are reported below.
No apparent proxy advantage may be represented as a verified
published-method comparison.

\begin{table}[H]
\centering\small
\caption{Experiment III: all 26 executed procedures, with descriptive macro-average normalized AMSE $Q_N$ across four signals and three root-SNR conditions for each $N$. A dagger denotes an implementation requiring audit; a double dagger denotes a provisional proxy and \emph{not} the exact original literature method. Numbers match the comparator table of the main article; citations identify the target methods. These results describe executed codes, not verified comparative performance of the original published algorithms.}
\label{tab:s5-e3-ranking}
\setlength{\tabcolsep}{5pt}
\begin{tabular}{r p{.58\linewidth}rrr}
\toprule
No. & Procedure and source & $N=512$ & $N=1024$ & $N=2048$ \\
\midrule
1 & Fixed SCOPE (normal; $r=0$) & 0.7100 & 0.7413 & 0.8004 \\
2 & Fixed N-SCOPE (normal; $r=1$) & 0.7368 & 0.8078 & 0.8855 \\
3 & SCOPE SURE (normal) & 0.6813 & 0.7003 & 0.7308 \\
4 & N-SCOPE SURE (normal) & 0.6100 & 0.6328 & 0.6599 \\
5 & BAMS$^{\dagger}$ \citep{VidakovicRuggeri2001} & --- & --- & --- \\
6 & EBayesThresh posterior median$^{\dagger}$ \citep{JohnstoneSilverman2005EBayes} & 0.7379 & 0.7499 & 0.7754 \\
7 & Raised-cosine Bayes$^{\dagger}$ \citep{ReinaSousa2026RaisedCosine} & 0.7550 & 0.8167 & 0.9181 \\
8 & Epanechnikov Bayes (1D)$^{\dagger}$ \citep{BarriosSousa2025Epanechnikov} & 1.0837 & 0.9884 & 0.9013 \\
9 & Smooth SCAD (stabilized)$^{\dagger}$ \citep{KulkarniEtAl2026SmoothSCAD} & 0.7572 & 0.7457 & 0.7616 \\
10 & SureShrink (hybrid)$^{\dagger}$ \citep{DonohoJohnstone1995} & 1.0195 & 1.1210 & 1.1954 \\
11 & FDR ($q=0.05$)$^{\dagger}$ \citep{AbramovichEtAl2006FDR} & 0.8759 & 0.8816 & 0.9049 \\
12 & Empirical Bayes block (1D)$^{\ddagger}$ \citep{AbramovichBesbeasSapatinas2002BlockEB} & 1.0550 & 1.2282 & 1.4430 \\
13 & BayesShrink (1D)$^{\dagger}$ \citep{ChangYuVetterli2000BayesShrink} & 1.0379 & 1.2439 & 1.4118 \\
14 & NeighCoeff$^{\dagger}$ \citep{CaiSilverman2001Neigh} & 0.9937 & 1.0552 & 1.0831 \\
15 & NeighShrink--SURE (1D)$^{\dagger}$ \citep{ZhouCheng2008NeighShrinkSURE} & 0.7685 & 0.8269 & 0.8549 \\
16 & SURE-LET (1D two-term)$^{\dagger}$ \citep{BluLuisier2007SURELET} & 0.8550 & 0.8695 & 0.8947 \\
17 & Sendur--Selesnick BiShrink (1D)$^{\dagger}$ \citep{SendurSelesnick2002BiShrink} & 0.7664 & 0.8763 & 1.0191 \\
18 & NeighBlock$^{\dagger}$ \citep{CaiSilverman2001Neigh} & 1.3614 & 1.3977 & 1.3404 \\
19 & SureBlock (hybrid)$^{\dagger}$ \citep{CaiZhou2009SureBlock} & 0.8027 & 0.8484 & 0.8669 \\
20 & BlockJS$^{\dagger}$ \citep{Cai1999BlockJS} & 1.5494 & 1.5962 & 1.7023 \\
21 & Jeffreys-prior EB (1D)$^{\dagger}$ \citep{FigueiredoNowak2001Jeffreys} & 0.8105 & 0.9557 & 1.2298 \\
22 & Local Bayesian Wiener (1D)$^{\dagger}$ \citep{MihcakEtAl1999LocalWiener} & 1.0029 & 1.2883 & 1.6887 \\
23 & Context-based shrinkage (1D)$^{\ddagger}$ \citep{ChangYuVetterli2000Context} & 1.1050 & 1.3465 & 1.5992 \\
24 & Hidden Markov tree (1D EM)$^{\ddagger}$ \citep{CrouseNowakBaraniuk1998HMT} & 0.6094 & 0.6022 & 0.6220 \\
25 & Universal hard threshold \citep{DonohoJohnstone1994} & 1.0000 & 1.0000 & 1.0000 \\
26 & Universal soft threshold \citep{DonohoJohnstone1994} & 3.1239 & 3.4744 & 3.7599 \\
\bottomrule
\end{tabular}
\end{table}


An aggregate can disguise signal-specific reversals. To document the
underlying conditions, the next three tables provide the four SCOPE
recipes plus three comparator codes, including one explicitly
identified HMT proxy, for \emph{every} signal and root-SNR setting at
all three sizes. These numerical comparisons do not use the test
signal for selection. The full 26-code condition-specific tables,
including replication-level MSEs and paired differences, are provided
with the machine-readable outputs.
% CONFIRM (authors): these files are in the public repository. If they hold the paired
% differences, the TODO table in the precision subsection can be filled from them directly.
BAMS estimates are withheld pending verification, consistently with the main article.

\begin{table}[H]
\centering\small
\caption{Experiment III, $N=512$: AMSE for each signal and root SNR; independent $M=1000$ paired replications. EBayes and NeighShrink-SURE require implementation audit; HMT is a provisional proxy.}
\label{tab:s5-e3-amse-512}
\setlength{\tabcolsep}{5pt}
\scriptsize
\setlength{\tabcolsep}{3pt}
\begin{tabular}{llrrrrrrr}
\toprule
Signal & SNR & Fixed S & Fixed N & S-SURE & N-SURE & EB med. & NeighS & HMT proxy \\
\midrule
Doppler & 3 & 0.2238 & 0.1832 & 0.2315 & 0.1802 & 0.2548 & 0.1823 & 0.2180 \\
Doppler & 5 & 0.2843 & 0.2158 & 0.2843 & 0.2007 & 0.3174 & 0.1965 & 0.2584 \\
Doppler & 7 & 0.3086 & 0.2281 & 0.3020 & 0.2079 & 0.3383 & 0.2060 & 0.2705 \\
HeaviSine & 3 & 0.0970 & 0.0951 & 0.0875 & 0.0845 & 0.0854 & 0.0955 & 0.1052 \\
HeaviSine & 5 & 0.1097 & 0.1127 & 0.1062 & 0.1090 & 0.1041 & 0.1197 & 0.1061 \\
HeaviSine & 7 & 0.1186 & 0.1240 & 0.1152 & 0.1189 & 0.1214 & 0.1355 & 0.1191 \\
Bumps & 3 & 0.4764 & 0.3814 & 0.4176 & 0.3783 & 0.4581 & 0.3916 & 0.3792 \\
Bumps & 5 & 0.5114 & 0.3689 & 0.4479 & 0.3635 & 0.4892 & 0.3808 & 0.3933 \\
Bumps & 7 & 0.5421 & 0.3570 & 0.4690 & 0.3532 & 0.5052 & 0.3705 & 0.3956 \\
Blocks & 3 & 0.2042 & 0.2797 & 0.2050 & 0.2050 & 0.2250 & 0.3020 & 0.1406 \\
Blocks & 5 & 0.1793 & 0.2751 & 0.1791 & 0.1791 & 0.1990 & 0.2970 & 0.1465 \\
Blocks & 7 & 0.1617 & 0.2628 & 0.1627 & 0.1627 & 0.1837 & 0.2870 & 0.1208 \\
\bottomrule
\end{tabular}
\end{table}

\begin{table}[H]
\centering\small
\caption{Experiment III, $N=1024$: AMSE for each signal and root SNR; independent $M=1000$ paired replications. EBayes and NeighShrink-SURE require implementation audit; HMT is a provisional proxy.}
\label{tab:s5-e3-amse-1024}
\setlength{\tabcolsep}{5pt}
\scriptsize
\setlength{\tabcolsep}{3pt}
\begin{tabular}{llrrrrrrr}
\toprule
Signal & SNR & Fixed S & Fixed N & S-SURE & N-SURE & EB med. & NeighS & HMT proxy \\
\midrule
Doppler & 3 & 0.1356 & 0.1134 & 0.1399 & 0.1032 & 0.1441 & 0.1068 & 0.1189 \\
Doppler & 5 & 0.1618 & 0.1334 & 0.1648 & 0.1245 & 0.1779 & 0.1253 & 0.1447 \\
Doppler & 7 & 0.1804 & 0.1493 & 0.1821 & 0.1384 & 0.1974 & 0.1374 & 0.1639 \\
HeaviSine & 3 & 0.0632 & 0.0647 & 0.0503 & 0.0509 & 0.0511 & 0.0592 & 0.0607 \\
HeaviSine & 5 & 0.0699 & 0.0740 & 0.0572 & 0.0597 & 0.0606 & 0.0723 & 0.0598 \\
HeaviSine & 7 & 0.0780 & 0.0814 & 0.0688 & 0.0710 & 0.0724 & 0.0835 & 0.0730 \\
Bumps & 3 & 0.2826 & 0.2253 & 0.2740 & 0.2243 & 0.2965 & 0.2419 & 0.2385 \\
Bumps & 5 & 0.3127 & 0.2348 & 0.2911 & 0.2347 & 0.3168 & 0.2427 & 0.2413 \\
Bumps & 7 & 0.3193 & 0.2428 & 0.2946 & 0.2460 & 0.3239 & 0.2474 & 0.2589 \\
Blocks & 3 & 0.1594 & 0.2211 & 0.1583 & 0.1583 & 0.1662 & 0.2326 & 0.1012 \\
Blocks & 5 & 0.1251 & 0.2027 & 0.1220 & 0.1220 & 0.1275 & 0.2184 & 0.0779 \\
Blocks & 7 & 0.1054 & 0.1844 & 0.1037 & 0.1037 & 0.1184 & 0.2040 & 0.0718 \\
\bottomrule
\end{tabular}
\end{table}

\begin{table}[H]
\centering\small
\caption{Experiment III, $N=2048$: AMSE for each signal and root SNR; independent $M=1000$ paired replications. EBayes and NeighShrink-SURE require implementation audit; HMT is a provisional proxy.}
\label{tab:s5-e3-amse-2048}
\setlength{\tabcolsep}{5pt}
\scriptsize
\setlength{\tabcolsep}{3pt}
\begin{tabular}{llrrrrrrr}
\toprule
Signal & SNR & Fixed S & Fixed N & S-SURE & N-SURE & EB med. & NeighS & HMT proxy \\
\midrule
Doppler & 3 & 0.0815 & 0.0725 & 0.0811 & 0.0580 & 0.0804 & 0.0614 & 0.0755 \\
Doppler & 5 & 0.0932 & 0.0835 & 0.0927 & 0.0707 & 0.0972 & 0.0722 & 0.0821 \\
Doppler & 7 & 0.1062 & 0.0900 & 0.1081 & 0.0754 & 0.1156 & 0.0773 & 0.0998 \\
HeaviSine & 3 & 0.0439 & 0.0456 & 0.0274 & 0.0285 & 0.0302 & 0.0357 & 0.0318 \\
HeaviSine & 5 & 0.0506 & 0.0524 & 0.0366 & 0.0377 & 0.0401 & 0.0457 & 0.0395 \\
HeaviSine & 7 & 0.0555 & 0.0571 & 0.0457 & 0.0465 & 0.0490 & 0.0532 & 0.0515 \\
Bumps & 3 & 0.1690 & 0.1390 & 0.1700 & 0.1408 & 0.1802 & 0.1465 & 0.1395 \\
Bumps & 5 & 0.1854 & 0.1543 & 0.1850 & 0.1557 & 0.2007 & 0.1557 & 0.1560 \\
Bumps & 7 & 0.1913 & 0.1610 & 0.1898 & 0.1618 & 0.2101 & 0.1602 & 0.1711 \\
Blocks & 3 & 0.0932 & 0.1310 & 0.0927 & 0.0927 & 0.0953 & 0.1408 & 0.0558 \\
Blocks & 5 & 0.0806 & 0.1264 & 0.0802 & 0.0802 & 0.0805 & 0.1351 & 0.0432 \\
Blocks & 7 & 0.0702 & 0.1182 & 0.0654 & 0.0654 & 0.0713 & 0.1243 & 0.0403 \\
\bottomrule
\end{tabular}
\end{table}


\subsection{Precision, reproducibility, and comparator audit}
\label{sec:supp-precision}
The Monte Carlo standard error of each AMSE is the standard deviation
of its replication-level MSE divided by $\sqrt{1000}$.
Paired standard errors use the standard deviation of
\emph{within-replication} differences divided by $\sqrt{1000}$.
A conventional descriptive $95\%$ Monte Carlo interval for the
expected MSE difference is its sample mean plus or minus $1.96$
paired standard errors.
% TODO (authors): insert the paired-difference table here before submission.
% For N=1024 (and, if space permits, N=512 and 2048), each signal, and root SNR 3,5,7,
% report mean(L_A - L_B) and sd(L_A - L_B)/sqrt(1000) over the 1000 replications for
%   A = N-SCOPE SURE, B = NeighShrink--SURE (1D)   [closest comparator; Doppler gap 0.0008 at SNR 5]
%   A = N-SCOPE SURE, B = SCOPE SURE               [source of the ~10% aggregate gain]
% MATLAB, with LA, LB the 1000-by-1 vectors of replication MSEs:
%   d = LA - LB;  m = mean(d);  se = std(d)/sqrt(numel(d));  ci = m + [-1.96 1.96]*se;
% Suggested layout (Table S-paired):
%   Signal & SNR & NS-SURE - NeighShrink (SE) & NS-SURE - SCOPE SURE (SE) \\
These intervals quantify Monte Carlo
precision conditional on the simulation design, not robustness to
other signal families or noise models. Where figures average
percentage gains across SNR, the average is a descriptive summary
and no pooled confidence interval is inferred without the
corresponding replication-level calculation.

The comparison set includes algorithmic modifications. For example,
the SureShrink code includes a sparse-level hybrid fallback, the
Smooth SCAD implementation includes an author-specified stabilization,
and the LET comparator uses a one-dimensional orthonormal two-term rule
rather than the original image-oriented method. The hidden-Markov-tree,
context, and empirical-Bayes-block codes are explicitly flagged as
provisional proxies. The SURE-LET, SureBlock, and Jeffreys-prior codes
are modified one-dimensional implementations; together with the three
proxies, they form the group that the comparator table of the main
article describes as provisional or modified.  The primary conclusions
in Section~5 of the main article are about SCOPE's morphology-dependent
behavior and the value of non-oracle adaptation, not publication-level
claims of superiority over any unverified literature method.


\bibliographystyle{apalike}
\bibliography{references}


\end{document}
